\documentclass[../DoAn.tex]{subfiles}
\begin{document}

\section{Danh sách REST API của dịch vụ Java}
Dịch vụ Java cung cấp API cho xác thực, phòng, yêu cầu đặt phòng, bảo trì, cấu hình, quyền tạo mã và thông báo. Các thao tác cần phiên đăng nhập được kiểm tra bằng token Bearer; đăng ký và đăng nhập có quy tắc xác thực riêng. Bảng~\ref{tab:api} nêu các nhóm điểm cuối tiêu biểu, không phải danh sách đầy đủ của mã nguồn hiện tại.

\begin{table}[H]\centering\small\renewcommand{\arraystretch}{1.15}
\caption{Các nhóm REST API tiêu biểu của dịch vụ Java}\label{tab:api}
\begin{tabularx}{\textwidth}{|p{1.5cm}|p{5.5cm}|X|}\hline
\textbf{Phương thức} & \textbf{Đường dẫn} & \textbf{Chức năng} \\ \hline
POST & \texttt{/api/auth/login} & Đăng nhập, tạo phiên \\ \hline
GET & \texttt{/api/rooms} & Đọc danh sách phòng \\ \hline
GET & \texttt{/api/configuration} & Đọc quy tắc nghiệp vụ \\ \hline
GET & \texttt{/api/maintenance} & Đọc lịch bảo trì \\ \hline
GET, POST & \texttt{/api/bookings} & Đọc và tạo yêu cầu \\ \hline
GET & \texttt{/api/room-occupancy} & Đọc khung giờ phòng đang chờ duyệt hoặc đã duyệt \\ \hline
POST & \texttt{/api/bookings/}\newline\texttt{\{id\}/review} & Duyệt hoặc từ chối \\ \hline
POST & \texttt{/api/bookings/}\newline\texttt{\{id\}/temporary-pin} & Lưu mã tạm thời và thông tin giao lệnh \\ \hline
GET & \texttt{/api/pin-permissions} & Xem quyền theo phòng \\ \hline
GET & \texttt{/api/notifications} & Đọc thông báo \\ \hline
\end{tabularx}\end{table}

Ứng dụng hiện giao tiếp OneIoT trực tiếp qua module Android. Các điểm cuối HTTP của gateway Python là công cụ thử nghiệm độc lập, không nằm trong đường truyền chính của ứng dụng.

\section{Hướng dẫn triển khai và kiểm tra}
\paragraph{Bản cục bộ.} Có thể cấu hình ứng dụng để lưu dữ liệu trên một thiết bị, phù hợp cho trình diễn không cần máy chủ. Bản phát hành Android sau khi đóng gói có thể chạy độc lập với công cụ phát triển. Tài khoản trình diễn được khởi tạo sẵn trong dữ liệu mẫu.

\paragraph{Bản dùng backend.} Dịch vụ Java sử dụng Spring Boot, Spring Data JPA và cơ sở dữ liệu H2 ở chế độ tệp theo cấu hình mặc định. Có thể điều chỉnh thông tin kết nối cơ sở dữ liệu theo môi trường triển khai. Để ứng dụng trao đổi dữ liệu, dịch vụ cần đang chạy và địa chỉ truy cập phải dùng được từ điện thoại. Kết nối công khai hiện dùng tunnel thử nghiệm, nên cần kiểm tra trạng thái đường truyền trước khi trình diễn.

\paragraph{Phát triển và kiểm thử.} Ứng dụng di động được xây dựng bằng React Native và công cụ Android; các kiểm thử tự động tập trung vào quy tắc nghiệp vụ và có kiểm tra định dạng, kiểu dữ liệu. Dịch vụ Java có các kiểm thử tích hợp với cơ sở dữ liệu thử nghiệm. Khi kiểm tra chế độ máy chủ, cần xác nhận đăng nhập, đọc/ghi dữ liệu và kết nối từ điện thoại.

\paragraph{Kết nối khóa số.} Gắn khóa đang quản lý với một phòng khóa số trong ứng dụng. Người dùng xác thực OneIoT trong phiên kết nối; thông tin xác thực không được lưu trong kho mã nguồn. Nếu ứng dụng chưa kết nối, mã mới tạo có thể ở trạng thái chờ gửi và được thử gửi lại khi kết nối được thiết lập, với điều kiện mã còn hiệu lực và phòng vẫn gắn với khóa đó.

\section{Biên bản kiểm tra trên thiết bị}
Bảng~\ref{tab:bienban} tóm tắt phạm vi minh chứng trong biên bản thiết bị của ứng dụng. Kết quả được ghi cho chế độ cục bộ; kiểm thử khóa vật lý và hệ thống nhiều thiết bị là hạng mục riêng.

\begin{table}[H]\centering\small\renewcommand{\arraystretch}{1.15}
\caption{Phạm vi kiểm tra bản Release trên điện thoại Android}\label{tab:bienban}
\begin{tabularx}{\textwidth}{|X|p{3.0cm}|}\hline
\textbf{Hạng mục} & \textbf{Ghi nhận} \\ \hline
Cài và khởi động bản Release không cần Metro & Đã kiểm tra \\ \hline
Đăng nhập tài khoản mẫu, điều hướng theo vai trò & Đã kiểm tra \\ \hline
Danh sách phòng theo tầng và chức năng đặt phòng cục bộ & Đã kiểm tra \\ \hline
Kiểm thử tự động ứng dụng ở lần rà soát gần nhất & 39/39 ca đạt \\ \hline
Kiểm thử khóa thật áp dụng mã và đồng bộ nhiều điện thoại & Chưa có xác nhận đầu cuối \\ \hline
\end{tabularx}\end{table}
\end{document}
